\scanpage{627}
\mspage{15}
\[
TC\subset C\tag*{(2.7)}
\]
de $\mathrm{Fl}\,M$, satisfaisant les conditions

% slanted
\margin{$C,TC$ ``spéciales''}
\[
\left|\begin{aligned}
&\left\{\begin{aligned}
\text{a)}\quad &\deleted{\unsure{\Phi}}C,TC\text{ sont }Q\text{-closes à gauche}\\
\text{b)}\quad &\deleted{\unsure{\Phi}}C,TC\text{ sont }Q\text{-engendrées à g. par}\\
&\text{des \uline{petites} parties } C_0,TC_0\text{ de }C,TC\\
&\text{respectivement}
\end{aligned}\right.\\
&\text{c)}\quad TC\subset W,\quad TF\overset{\mathrm{def}}{=}C_*\subset W\\
&\text{d)}\quad\text{On a l'une des \inserted{trois} conditions équivalentes}\\
&\qquad\text{supplémentaires de (2.4), i.e}\\
&\qquad\alpha)\ TC=C\cap W\quad\text{i.e}\quad C\cap W\subset TC\\
&\qquad\beta)\ TF=F\cap W\quad\text{i.e.}\quad F\cap W\subset TF\\
&\qquad\gamma)\ W=TF\circ TC.
\end{aligned}\right.\tag*{(2.8)}
\]

\uline{Proposition} 2.1. Soit

\diagram{\begin{tikzcd}[ampersand replacement=\&,column sep=large]
\text{(2.1.1)} \& M'\arrow[r,shift right,"u^*"'] \& M\arrow[l,shift right,"u_!"']
\end{tikzcd}}

un couple de foncteurs adjoints, $M$ étant
muni d'une structure de Quillen close
$(W,C,F)$. Posons

\[
\left\{\begin{aligned}
&\deleted{C'=u_!(C)^{\sim}}\\
&\deleted{TC'=u_!(TC)^{\sim}}\\
&\deleted{F'\unsure{.}}
\end{aligned}\right.\quad
\left\{\begin{aligned}
&W'=(u^*)^{-1}(W)\\
&F'=(u^*)^{-1}(W)\\
&TF'=F'\cap W'=(u^*)^{-1}(TF)\\
&\deleted{\unsure{T}}C'=u_!(C)^{\sim}\\
&TC'=u_!(TC)^{\sim}.
\end{aligned}\right.\tag*{(2.1.2)}
\]

$1^\circ)$ On a alors

\[
\left\{\begin{aligned}
&C'\longleftrightarrow TF',\qquad TC'\longleftrightarrow F'\\
&T\unsure{\cdot} C'\subset C',\qquad TF'\subset F',
\end{aligned}\right.\tag*{(2.1.3)}
\]

Pour que les données (2.1.\unsure{1|2}) dans $M'$ \unsure{//|))} correspondent à un triple de Quillen $(W',C',F')$, il
faut et il suffit que les conditions \unsure{suivantes|équiv. suivantes}
soient satisfaites :
% archive sheet number: 626
